\section{Analisis Kelemahan Cipher Klasik}

\subsection{Frequency Analysis}

Dalam bahasa Inggris, frekuensi huruf kira-kira: E (12\%), T (9\%), A (8\%), ... Dengan ciphertext panjang, huruf yang paling sering dalam ciphertext kemungkinan memetakan ke E, T, atau A. Pendekatan ini efektif untuk Caesar dan substitusi monoalfabetik.

\subsection{Kasiski Examination}

Untuk Vigenère, jika plaintext memiliki pola berulang (misal THE) dan kunci pendek, pola tersebut akan muncul dalam ciphertext dengan jarak kelipatan panjang kunci. Kasiski examination \cite{stallings2017}:
\begin{enumerate}
  \item Cari n-gram berulang (2--4 huruf) dalam ciphertext.
  \item Ukur jarak antar kemunculan; jarak tersebut sering merupakan kelipatan panjang kunci.
  \item Ambil $\gcd$ dari beberapa jarak untuk menduga panjang kunci $m$.
  \item Pisahkan ciphertext menjadi $m$ aliran (posisi $i$, $i+m$, $i+2m$, \ldots); setiap aliran bersifat seperti Caesar.
  \item Lakukan frequency analysis pada setiap aliran untuk memulihkan huruf kunci.
\end{enumerate}

\textbf{Index of Coincidence} dapat digunakan untuk memperkirakan panjang kunci: jika $m$ benar, index of coincidence tiap aliran mendekati nilai bahasa asli; jika salah, mendekati nilai acak.

\subsection{Ringkasan Serangan Spesifik}

\begin{itemize}
  \item \textbf{Caesar}: brute force (25 kunci) dan frequency analysis.
  \item \textbf{Monoalfabetik}: frequency analysis dan pola digraph/trigraph.
  \item \textbf{Affine}: brute force pada pasangan $(a,b)$, karena $a$ hanya 12 nilai valid.
  \item \textbf{Vigenere}: Kasiski dan Index of Coincidence untuk menebak panjang kunci.
  \item \textbf{Playfair}: analisis frekuensi digraph, terutama jika ciphertext panjang.
  \item \textbf{Hill}: known-plaintext attack, karena sistem linier dapat diselesaikan dengan aljabar linear.
  \item \textbf{Transposisi}: anagramming dan analisis struktur, terutama bila format pesan diketahui.
\end{itemize}
\cite{stallings2017,menezes1996}

\begin{konsep}
Cipher klasik tidak aman untuk penggunaan modern. Mereka diajarkan untuk memahami evolusi kriptografi dan dasar kriptanalisis. Algoritma modern (AES, RSA) dirancang untuk menahan serangan yang canggih.
\end{konsep}
